\chapter{Контрольная по занятиям №1--2}\label{mock:start}
\textbf{90 минут, 100 баллов.} Вариант охватывает лекции и семинары №1--2. Работайте без калькулятора, программ, учебных материалов и ответов. Все вычисления выполняйте на бумаге; показывайте существенные промежуточные шаги. Время решения --- примерно 80 минут, проверки --- 10 минут. Распределение баллов предназначено для самопроверки тренировочного варианта.

Во всех заданиях с машинными словами ширина 8 битов, если явно не указано расширение до 16. Знаковая интерпретация --- дополнительный код; для смещённого кода bias $=128$. Биты нумеруются справа с нуля. Флаги записывайте в порядке CF, OF, ZF, SF; для вычитания используются правила SUB x86. Подпункты независимы.

\section{К1. Переводы (10 баллов, 10 минут)}\label{task:k1}
\begin{enumerate}
    \item Переведите $173_{10}$ в двоичную, восьмеричную и шестнадцатеричную системы. Покажите способ перевода. \textbf{(5 баллов)}
    \item Переведите $2E6_{16}$ в двоичную и восьмеричную системы группировкой битов. \textbf{(5 баллов)}
\end{enumerate}

\section{К2. Арифметика (10 баллов, 8 минут)}\label{task:k2}
Разрядность в этом задании не ограничена.
\begin{enumerate}
    \item Вычислите $3A6_{16}+7B_{16}$ с указанием переносов и проверьте ответ в десятичной системе. \textbf{(5 баллов)}
    \item Найдите частное и остаток от $264_8:7_8$. Выполните проверку $A=BQ+R$. \textbf{(5 баллов)}
\end{enumerate}

\section{К3. Диапазоны (10 баллов, 8 минут)}\label{task:k3}
\begin{enumerate}
    \item Укажите диапазон и количество различных значений 10-битного дополнительного кода. \textbf{(4 балла)}
    \item Сколько битов нужно unsigned-числу со значениями от 0 до 1500? Обоснуйте минимальность. \textbf{(3 балла)}
    \item Сколько битов нужно для диапазона $[-600,600]$ в дополнительном коде? Обоснуйте минимальность. \textbf{(3 балла)}
\end{enumerate}

\section{К4. Коды (15 баллов, 12 минут)}\label{task:k4}
\begin{enumerate}
    \item Представьте $-53$ в 8-битных прямом, обратном, дополнительном кодах и коде со смещением 128. Запишите биты и hex-коды. \textbf{(8 баллов)}
    \item Прочитайте 8-битное $D9_{16}$ как unsigned и как дополнительный код. \textbf{(4 балла)}
    \item Выполните знаковое расширение $D9_{16}$ до 16 битов. \textbf{(3 балла)}
\end{enumerate}

\section{К5. Результаты и флаги (20 баллов, 14 минут)}\label{task:k5}
Выполните операции в 8 битах. Для каждой укажите hex-код результата, его знаковое значение и CF, OF, ZF, SF. Обоснуйте CF и OF. \textbf{По 5 баллов за подпункт.}
\begin{enumerate}
    \item $63+65$;
    \item $(-90)+(-50)$;
    \item $25-70$;
    \item $(-100)-40$.
\end{enumerate}

\section{К6. Побитовые операции (10 баллов, 10 минут)}\label{task:k6}
Даны 8-битные $x=CA_{16}$, $y=75_{16}$.
\begin{enumerate}
    \item Найдите $x$ AND $y$, $x$ OR $y$, $x$ XOR $y$, NOT $x$. Для каждого ответа запишите биты и hex-код. \textbf{(8 баллов)}
    \item Независимо замените младшую тетраду исходного $x$ на $3_{16}$, сохранив старшую. Укажите маску, операции и результат. \textbf{(2 балла)}
\end{enumerate}

\section{К7. Сдвиги (15 баллов, 10 минут)}\label{task:k7}
Исходное 8-битное слово $x=A9_{16}$. Каждый подпункт выполняйте над исходным словом. \textbf{По 3 балла за подпункт.}
\begin{enumerate}
    \item Логический сдвиг вправо на 2: биты, hex и unsigned-значение.
    \item Арифметический сдвиг вправо на 2: биты, hex и знаковое значение.
    \item Циклический сдвиг влево на 1: биты и hex-код.
    \item Знаковое расширение до 16 битов: биты и hex-код.
    \item Прочитайте исходное слово как знаковое число, разделите его на 4 с округлением к нулю и сравните с подпунктом 2.
\end{enumerate}

\section{К8. Проверка чужого решения (10 баллов, 8 минут)}\label{task:k8}
Студент складывает 8-битные коды $FD_{16}$ и $03_{16}$:

\begin{quote}
«Сумма равна $100_{16}$, её и сохраняем в 8-битном результате. CF=1, поэтому OF=1. Первый операнд отрицательный, значит SF=1. ZF=0, ведь полная сумма не нулевая».
\end{quote}

Укажите сохранённый 8-битный код, точную сумму знаковых значений и правильные четыре флага. Исправьте рассуждения студента.

\section*{Последние 10 минут: проверка}
\begin{itemize}
    \item У каждого числа указан нужный формат или основание?
    \item В ответе сохранено требуемое количество битов?
    \item CF и OF проверены отдельно? У вычитания CF определён по займу?
    \item SF и ZF вычислены по сохранённому результату?
    \item Все сдвиги и маски применены к нужному исходному слову?
    \item Остаток неотрицателен и меньше делителя?
\end{itemize}
Закончив работу, переходите к \hyperref[answers:mock]{кратким ответам} и \hyperref[solutions:mock]{решениям с разбалловкой}.
